\documentclass[12pt]{article}
\usepackage{ficha}

\begin{document}

\ficha{Função Quadrática, Módulo, Funções Definidas por Ramos e Função Raiz}

\instrucoes{Resolve os exercícios, apresentando os cálculos. Apresenta os conjuntos-solução em notação de intervalos.}

\secao{Função Quadrática}{Para cada função, determina o vértice, o eixo de simetria, os zeros, a concavidade e o contradomínio.}
\begin{exercicios}
  \item $f(x) = x^2 - 4x + 3$
  \item $f(x) = -x^2 + 6x - 5$
  \item $f(x) = 2x^2 + 4x - 6$
  \item $f(x) = (x - 3)^2 - 4$
  \item $f(x) = -2(x + 1)^2 + 8$
  \item $f(x) = x^2 + 2x + 5$
  \item Escreve $f(x) = x^2 - 6x + 7$ na forma $a(x - h)^2 + k$.
  \item Determina a expressão da parábola com vértice $(1, -2)$ que passa no ponto $(3, 6)$.
  \item Determina a expressão da parábola com zeros $-1$ e $4$ que passa no ponto $(0, -8)$.
\end{exercicios}

\secao{Equações e Inequações com Módulo}{Resolve, em $\R$.}
\begin{exercicios}
  \item $|x| = 5$
  \item $|x - 3| = 2$
  \item $|2x + 1| = 7$
  \item $|x + 4| = -1$
  \item $|x| < 3$
  \item $|x - 2| \le 5$
  \item $|x + 1| > 4$
  \item $|3x - 6| \ge 9$
  \item $|x - 1| = |x + 3|$
  \item $|x^2 - 5| = 4$
\end{exercicios}

\secao{Funções com Módulo}{Escreve cada função sem módulo (definida por ramos) e esboça o gráfico.}
\begin{exercicios}
  \item $f(x) = |x - 2|$
  \item $f(x) = |2x + 4|$
  \item $f(x) = |x| - 3$
  \item $f(x) = 2|x - 1| + 1$
  \item $f(x) = |x^2 - 4|$
  \item Para $f(x) = -|x + 2| + 3$, indica o máximo, os zeros e o contradomínio.
\end{exercicios}

\secao{Funções Definidas por Ramos}{Considera $f(x) = \begin{cases} x + 3 & \text{se } x < -1 \\ x^2 & \text{se } -1 \le x \le 2 \\ 6 - x & \text{se } x > 2 \end{cases}$}
\begin{exercicios}
  \item Calcula $f(-4)$, $f(-1)$, $f(0)$, $f(2)$ e $f(5)$.
  \item Esboça o gráfico de $f$.
  \item Indica o contradomínio de $f$.
  \item Determina os zeros de $f$.
  \item Resolve a equação $f(x) = 1$.
  \item Resolve a inequação $f(x) \ge 2$.
  \item Indica os intervalos de monotonia de $f$.
\end{exercicios}

\secao{Função Raiz Quadrada e Raiz Cúbica}{Resolve.}
\begin{exercicios}
  \item Indica o domínio de $f(x) = \sqrt{x - 4}$.
  \item Indica o domínio de $g(x) = \sqrt{6 - 2x}$.
  \item Indica o domínio de $h(x) = \sqrt{x^2 - 9}$.
  \item Indica o domínio de $i(x) = \sqrt[3]{x - 1}$.
  \item Esboça o gráfico de $f(x) = \sqrt{x + 1} - 2$ e indica o domínio, o contradomínio e o zero.
  \item Resolve: $\sqrt{x + 5} = 3$
  \item Resolve: $\sqrt{2x - 1} = x - 2$
  \item Resolve: $\sqrt{x + 7} = x + 1$
  \item Resolve: $\sqrt[3]{x - 2} = -2$
\end{exercicios}

\secao{Desafios}{Resolve os problemas.}
\begin{exercicios}
  \item Uma bola é lançada e a sua altura, em metros, ao fim de $t$ segundos é $h(t) = -5t^2 + 20t + 1$. Qual é a altura máxima? Quando atinge o solo? (Arredonda às centésimas.)
  \item Com 40\,m de rede, quer-se vedar um terreno retangular encostado a um muro (só três lados precisam de rede). Que dimensões maximizam a área?
  \item Resolve: $|x - 1| + |x - 3| = 4$.
  \item Determina $k$ de modo que a função definida por $f(x) = x^2$ se $x \le 1$ e $f(x) = 2x + k$ se $x > 1$ seja contínua (os ramos «colem») em $x = 1$.
  \item Resolve: $\sqrt{x} + \sqrt{x - 5} = 5$.
\end{exercicios}

\end{document}
